\section{整体灵敏度与鲁棒性分析}
\label{sec:overall-sensitivity}

\subsection{核数变化下的逐图稳定性}
逐图继承规则使每张图的有效工期随可用核数单调不增：目标核数$K$的方案取已评价的$1,\ldots,K$核计划和整图单核基准中的最短Makespan。因而相邻核数的退化计数均为0，表\ref{tab:adjacent-core-tail}给出100张图上的退化计数与最大工期比。低核方案以空核列表补齐，工期和官方事件结果保持不变。

仅统计变慢图数还不足以决定是否增核。记同图、同场景相邻核数的选中方案工期比为
$R_{i,r,K}=F_{r,i,K+1}/F_{r,i,K}$；$R_{i,r,K}>1.1$表示增加一核后工期延长超过10\%。表\ref{tab:adjacent-core-tail}同时给出这类退化的图数和100图中的最大工期比。A的$1\to2$使用单核辅助诊断，其余A转移及B、C全部属于正式实验矩阵。
\begin{table}[H]\centering
\caption{相邻核数的退化尾部：超过10\%的图数／最大工期比}
\label{tab:adjacent-core-tail}
\begin{tabularx}{\linewidth}{cYYY}\toprule
核数变化 & A & B & C\\\midrule
$1\to2$ & $0\,/\,1.000$ & $0\,/\,1.000$ & $0\,/\,1.000$\\
$2\to3$ & $0\,/\,1.000$ & $0\,/\,1.000$ & $0\,/\,1.000$\\
$3\to4$ & $0\,/\,1.000$ & $0\,/\,1.000$ & $0\,/\,1.000$\\
$4\to5$ & $0\,/\,1.000$ & $0\,/\,1.000$ & $0\,/\,1.000$\\\bottomrule
\end{tabularx}
\end{table}

原始候选中\Case{067}的A场景四核工期为18,473,155 cycle，五核候选升至63,959,635 cycle；按继承规则，五核终态沿用四核计划并以空核列表补齐，因此最终五核工期仍为18,473,155 cycle。该例说明继承规则直接针对已评价计划的尾部退化，工期比较仍以完整官方事件模拟为准。

场景间仍按同一图、同一核数比较，但A/B/C的继承来源可能不同。固定五核计划的B/C同计划配对继续用于分离Cache硬件和调度适配效应；继承后的总体曲线则回答每张图在不超过$K$个核心时可取得的最短已评价工期。两种结果均回到逐图工期和事件时序解释，避免把场景差异混同为核数收益。

\subsection{逐图继承方案的来源与收益}
对每张图和场景，目标核数$K$的终态由$1,\ldots,K$核已评价计划及整图单核基准共同确定；若最短方案来自较低核数，则在目标配置中保留其计划并补齐空核列表。表\ref{tab:best-core-choice}统计五核终态的来源核数，整图基准单列；这组计数说明有多少图通过继承避免了原始五核候选的退化。五核最终结果已经是各图可用候选中的最短工期，因此不再把它与自身重复比较。
\begin{table}[H]\centering\small
\caption{100图五核终态的继承来源与最终速度比}\label{tab:best-core-choice}
\begin{tabular}{lrrrrrrr}\toprule
场景 & 整图基准 & 1核 & 2核 & 3核 & 4核 & 5核 & 五核平均加速比\\\midrule
A & 14 & 3 & 3 & 9 & 10 & 61 & 3.4132\\
B & 6 & 6 & 5 & 10 & 8 & 65 & 3.3862\\
C & 6 & 6 & 5 & 8 & 6 & 69 & 3.4272\\\bottomrule
\end{tabular}\end{table}
\samplefigure{core_choice_tail}{已评价计划中逐图选核相对五核计划的工期降幅分布}{fig:core-choice-tail}

五核来源仍以五核候选为主，但A、B、C分别有14、6、6张图直接回到整图单核基准，另有25、29、25张图继承了1--4核方案。由于这些继承方案在目标核数中补齐空核列表，四至五核的100图配对分别为A的61图缩短、39图不变，B的65图缩短、35图不变，C的69图缩短、31图不变，三种场景均无工期反升。\Case{067}是继承规则的典型例子：原始A五核候选为63,959,635 cycle，四核候选为18,473,155 cycle，最终五核沿用后者。

继承后的五核终态与整图基准逐图比较，A、B、C的慢例数均为0；五核平均加速比分别为3.4132、3.3862和3.4272，实例重采样95\%区间分别为$[3.1025,3.7151]$、$[3.0686,3.6824]$和$[3.1116,3.7245]$。这些区间是在继承后的逐图比值上重采样得到，反映100张图组成改变时均值的波动。

继承规则的收益集中在原始五核候选发生退化的图：先比较已评价工期，再决定是否沿用较低核计划；整图基准也作为合法候选保留。最终均值不再由退化尾部拉低，具体图的搬运、关键路径和容量恢复仍由所继承计划的官方事件记录解释。

表中结果是已评价计划集合的后验最短值；逐图取最小不把候选搜索时间计入工期指标。若运行条件允许每张图使用不超过$K$个核心，可直接采用继承后的计划；若必须预先固定一种核数，则应报告固定核数的原始候选结果，并将继承方案作为逐图调度策略单独比较。新增核数本身不构成改进，真正比较的是完整执行时间线。

\subsection{样本组成的重采样检验}
对固定的100张图有放回抽样1500次，每次重新计算继承后逐图速度比的均值，并取两端经验分位。五核A、B、C均值分别为3.4132、3.3862、3.4272，对应的95\%实例重采样区间分别为$[3.1025,3.7151]$、$[3.0686,3.6824]$和$[3.1116,3.7245]$。这些区间刻画图集合组成对均值的影响；逐图调度选择依据具体图和可用核数的工期。固定五核计划的B/C对数工期比仍采用同计划配对数据：$\log(F_B/F_C(X_B))$的均值为0.01159774，区间为$[0.00534581,0.01957431]$，C内适配比$\log(F_C(X_B)/F_C(X_C))$的均值为0.00112141，区间为$[0.00012292,0.00269912]$。

\subsection{评价预算灵敏度}
为考察候选评价次数对方案质量和计算开销的影响，另选取\Case{064}、\Case{070}、\Case{093}、\Case{100}、\Case{096}和\Case{055}六张代表图，在五核配置下固定候选生成顺序，逐步把累计完整评价预算设为2、5、10和20次。必评锚点先行，其余候选按轻量评分进入评价队列；表中成功率按该预算内已完成官方评价的尝试数计算，累计时间为官方评价耗时之和。该实验用于观察预算趋势，主矩阵的逐图结果仍按各自实验清单统计。
\begin{table}[H]\centering\scriptsize
\caption{六张代表图在不同评价预算下的平均结果}
\label{tab:budget-sensitivity}
\begin{tabularx}{\linewidth}{lYYYY}
\toprule
场景与指标 & 2次 & 5次 & 10次 & 20次\\\midrule
A 平均加速比 & 2.9283 & 3.2697 & 3.2697 & 3.4132\\
A 成功率 & 100.0\% & 83.3\% & 75.0\% & 71.6\%\\
A 累计评价时间/s & 0.675 & 1.211 & 1.943 & 3.969\\
B 平均加速比 & 2.9283 & 3.0916 & 3.0916 & 3.4132\\
B 成功率 & 100.0\% & 80.0\% & 65.0\% & 60.3\%\\
B 累计评价时间/s & 0.358 & 1.051 & 2.149 & 4.323\\
C 平均加速比 & 3.4166 & 3.4224 & 3.4308 & 3.4313\\
C 成功率 & 100.0\% & 100.0\% & 100.0\% & 100.0\%\\
C 累计评价时间/s & 0.363 & 0.911 & 1.812 & 3.390\\
\bottomrule
\end{tabularx}
\end{table}

表\ref{tab:budget-sensitivity}显示，A场景由2次增加到20次时，六图平均加速比由2.9283升至3.4132，主要增益出现在前5次评价和少数候选较多的图；C场景从3.4166升至3.4313，增益较小而累计评价时间近似随预算增加。B场景的平均值在5至10次之间保持不变，20次时达到3.4132，但部分候选因依赖环未通过，成功率随预算下降。因而增加预算能够改善部分图的已评最优工期，却同时增加失败尝试和评价成本；对超大图应先保留整图、继承方案与首个合法切块，再根据剩余时间追加评价。

\subsection{工期容差与搬运量}
问题三的Cache容量与带宽扰动已在第\ref{sec:question-three}节逐图分析。本节进一步考察选择准则本身：在已成功评价的候选中，允许工期不超过组内最短工期的1.005倍，正式1400组里有24组可改选搬运更少的方案，合计节省42.350 MB；容差放宽到1.01倍时，有35组可改选，合计节省78.277 MB，其中$1\,\mathrm{MB}=10^6\,\mathrm B$。这些改选都来自原先已评估的方案，说明小幅工期容差能够换取搬运节省。容差方案与原方案满足同一组可行性检查；实际采用哪一档，取决于任务对完成时间和搬运量的优先顺序。
